\section*{Capítulo \ref{ch:g4:geometry} --- Rectas y polígonos}

\begin{solution}{exo:g4:geometry:1}
Construcciones con el \cref{met:g4:geometry:draw}; el cuadradito
marca el \omterm{def:g3:shapes:rightangle}{ángulo recto} entre $d$ y la perpendicular.
\end{solution}

\begin{solution}{exo:g4:geometry:2}
E: los tres trazos horizontales son paralelos y cada uno es perpendicular
al vertical. H: los dos trazos verticales son paralelos y el
travesaño es perpendicular a los dos. N: los dos trazos verticales son
paralelos; la diagonal no es ninguna de las dos cosas. T: los dos trazos son
\omterm{def:g4:geometry:def}{perpendiculares}. Z: los trazos de arriba y de abajo son paralelos; la
diagonal no es ninguna de las dos cosas.
\end{solution}

\begin{solution}{exo:g4:geometry:3}
\omterm{def:g4:geometry:polygon}{Hexágono}; \omterm{def:g4:geometry:polygon}{cuadrilátero}; triángulo; \omterm{def:g4:geometry:polygon}{octógono}; \omterm{def:g4:geometry:polygon}{pentágono}. Un \omterm{def:g4:geometry:polygon}{polígono} tiene
tantos \omterm{def:g4:geometry:polygon}{vértices} como lados: $6$, $4$, $3$, $8$, $5$.
\end{solution}

\begin{solution}{exo:g4:geometry:4}
Un \omterm{def:g4:geometry:polygon}{pentágono} tiene $5$ lados. Desde $A$, los \omterm{def:g4:geometry:polygon}{vértices} no vecinos son $C$ y
$D$: dos diagonales, $[AC]$ y $[AD]$.
\end{solution}

\begin{solution}{exo:g4:geometry:5}
$OA = 4$ cm (el \omterm{def:g3:shapes:shapes}{radio}); $OM = 4$ cm (como todo punto de la circunferencia);
$AB = 8$ cm (un \omterm{def:g4:geometry:circle}{diámetro} mide el doble que el \omterm{def:g3:shapes:shapes}{radio}).
\end{solution}

\begin{solution}{exo:g4:geometry:6}
Abre el compás sobre un lado y lleva esa abertura a los otros
dos: si la abertura encaja exactamente en un segundo lado,
esos dos lados son iguales.
\end{solution}

\begin{solution}{exo:g4:geometry:7}
El recorrido se cierra en un \omterm{def:g4:geometry:polygon}{polígono} de $5$ lados --- un \omterm{def:g4:geometry:polygon}{pentágono} (¡no
regular!).
\end{solution}

\begin{solution}{exo:g4:geometry:8}
«Dos \omterm{def:g4:geometry:def}{rectas perpendiculares} siempre se cortan»: verdadero --- se cortan justo en el
\omterm{def:g3:shapes:rightangle}{ángulo recto}.

«Dos \omterm{def:g4:geometry:def}{rectas paralelas} no se cortan nunca»: verdadero --- esa es su definición.

«Una recta puede ser paralela a una recta y perpendicular a
otra»: verdadero --- en papel cuadriculado, una recta horizontal es paralela a
todas las horizontales y perpendicular a todas las verticales.
\end{solution}

\begin{solution}{exo:g4:geometry:9}
Las rectas $d$ y $f$ parecen \omterm{def:g4:geometry:def}{paralelas} --- y siempre lo son: dos
\omterm{def:g4:geometry:def}{rectas perpendiculares} a una misma recta son \omterm{def:g4:geometry:def}{paralelas}
(la \cref{prop:g6:lines:facts} lo hace oficial).
\end{solution}

\begin{solution}{exo:g4:geometry:10}
Un programa correcto: «1. Traza un cuadrado $ABCD$ de lado $4$ cm.
2. Marca el punto $O$ donde se cortan sus diagonales. 3. Traza la circunferencia
de centro $O$ y \omterm{def:g3:shapes:shapes}{radio} $2$ cm.» La circunferencia toca cada lado en
su punto medio.
\end{solution}

\begin{solution}{exo:g4:geometry:11}
\omterm{def:g4:geometry:polygon}{Cuadrilátero}: $2$ diagonales. \omterm{def:g4:geometry:polygon}{Pentágono}: $5$. \omterm{def:g4:geometry:polygon}{Hexágono}: $9$. Cada \omterm{def:g4:geometry:polygon}{vértice}
nuevo añade cada vez más diagonales --- los recuentos crecen cada vez
más deprisa ($+3$, luego $+4$, \dots).
\end{solution}
